%=======================================================================
\section{Dataset and Preprocessing}

\subsection{Dataset}
The dataset is organised as one folder per class (\texttt{images\_by\_class/<ClassName>/}). After
loading and an 80/20 stratified split (random seed 42 for full reproducibility), the working set is:

\begin{center}
\begin{tabularx}{\textwidth}{@{}lY@{}}
\toprule
\thd{Property} & \thd{Value} \\
\midrule
Classes (\NClasses) & Glioma \textperiodcentered{} Meningioma \textperiodcentered{} Pituitary Tumor \\
Training images & \NTrain \\
Test images & \NTest \\
Total & \NTotal \\
Image size after preprocessing & $128\times128$, single-channel (grayscale) \\
Split & 80\,\% train / 20\,\% test, stratified by class, seed $=42$ \\
\bottomrule
\end{tabularx}
\end{center}

The split is \textbf{stratified} so that each class keeps the same proportion in train and test, and
the seed is fixed so the experiment can be repeated exactly.
Figure~\ref{fig:dataset-overview} summarises class counts and sample appearance.

\begin{figure}[htbp]
  \centering
  \includegraphics[width=\textwidth]{dataset}
  \caption{Dataset overview. \textbf{Left:} stratified train/test counts per class.
    \textbf{Right:} mosaic of training slices (raw resize preview; full preprocessing in
    Figure~\ref{fig:preprocessing-pipeline}).}
  \label{fig:dataset-overview}
\end{figure}

\subsection{Preprocessing pipeline}
Every image passes through an identical normalisation chain before any feature is computed. This
removes scanner- and acquisition-related nuisance variation so the descriptors describe anatomy, not
artefacts (Figure~\ref{fig:preprocessing-pipeline}):

\begin{enumerate}[leftmargin=1.6em,itemsep=2pt]
  \item \textbf{Grayscale conversion} --- colour channels (if any) are collapsed; MRI intensity is
        what matters.
  \item \textbf{Black-border cropping} --- the surrounding black region is detected (intensity
        $>0.05$) and cropped away so the brain fills the frame consistently.
  \item \textbf{Resize to $128\times128$} with anti-aliasing --- gives every image a common, fixed
        resolution so feature vectors are comparable.
  \item \textbf{Gaussian smoothing} ($\sigma=0.5$) --- suppresses high-frequency acquisition noise
        without blurring real structure.
  \item \textbf{Min--max normalisation to $[0,1]$} --- puts every image on the same intensity scale.
  \item \textbf{CLAHE} (Contrast-Limited Adaptive Histogram Equalisation, clip $=0.02$) --- locally
        boosts contrast so subtle tumor texture becomes visible to the descriptors.
\end{enumerate}

\begin{figure}[htbp]
  \centering
  \includegraphics[width=0.92\textwidth]{preprocessing_pipeline}
  \caption{Preprocessing pipeline. Each slice follows the same deterministic chain before
    descriptor extraction.}
  \label{fig:preprocessing-pipeline}
\end{figure}

\begin{notebox}[Why it matters:]
Cropping $+$ resizing make features spatially comparable; smoothing $+$ CLAHE make texture
descriptors (GLCM, LBP) respond to real tissue patterns instead of noise and global brightness. The
whole chain is cached to disk so re-runs are fast.
\end{notebox}


